% ============================================================
% 3.1-3.4 METHODOLOGY (core) -- Writer: Pablo Torres
% Target: about 1,000 words + equations (1)-(4).
% Brief: "Methodology or Approach: methods applied, implementation choices, and
% rationale." Template: "Each methodological decision should have a reason.
% Separate established methods from modifications or original decisions made by
% the team." Introduce every equation in the text and explain every symbol and
% unit. Results do NOT go here (they go in Section 4).
% The equations below follow the code exactly (src/progress.py, mixed_model.py,
% ministry.py, normal_check.py, final_figures.py). Keep the notation, explain
% it in your own words, and change it only together with Ignacio and Adam so
% that the whole report uses the same symbols.
% Start from outline entries 4, 4a, 4b, 4c, 5, 5a.
% ============================================================

We organised the analysis around three linked questions: how much grade variation belongs to the field, whether a university's particular field pattern persists, and how those differences affect access to a pooled ranking. Figure~\ref{fig:overview} maps this design. Reading left to right, we move from each dataset to the analysis it supports and then to the question and results subsection where we answer it. The arrows connect analytical tasks, not individual records across countries. We first define deviations and estimate variance components, then examine intake, simulate selection and describe uncertainty before presenting the prediction and ranking procedures in Sections~\ref{sec:methods-prediction}--\ref{sec:methods-computing}.
\subsection{Field Deviations and Variance Decomposition}
\label{sec:methods-decomposition}

We made fields comparable within a university and year by subtracting that university's average across its reported fields. In equation~\eqref{eq:deviation}, $\bar g_{ufy}$ is the mean final grade for university $u$, field $f$ and academic year $y$, measured in grade points on the 5--10 scale. The set $F_{uy}$ contains the fields reported for that university-year, and $|F_{uy}|$ is its size; $f'$ indexes the fields being averaged. The resulting deviation $d_{ufy}$ is positive when the field grades above its university-year average and negative when it grades below. We gave reported fields equal weight here, keeping the decomposition distinct from graduate-weighted selection.
\begin{equation}
  d_{ufy} \;=\; \bar g_{ufy} \;-\; \frac{1}{|F_{uy}|}\sum_{f' \in F_{uy}} \bar g_{uf'y}.
  \label{eq:deviation}
\end{equation}
We fitted nested least-squares models to separate a common field pattern from university-specific differences. In equation~\eqref{eq:decomposition}, $\alpha_{uy}$ captures the university-year level, $\beta_f$ the shared field effect, $\gamma_{uf}$ the additional effect of that field at that university, and $\varepsilon_{ufy}$ the unexplained annual variation. We first included $\alpha_{uy}$, then added $\beta_f$, and finally $\gamma_{uf}$. We measured each contribution by the reduction in residual sum of squares divided by the variation remaining after the university-year effects; the final residual forms the noise share. Thus the percentages describe differences within university-years, not all variation in students' grades. We used an $F$-test to compare models with and without the interaction across 1,528 observed university-field cells. We sum-coded field effects, imposing $\sum_f\beta_f=0$, so an effect measures distance from the average field rather than from a chosen reference subject.
\begin{equation}
  \bar g_{ufy} \;=\; \alpha_{uy} + \beta_f + \gamma_{uf} + \varepsilon_{ufy},
  \qquad \textstyle\sum_f \beta_f = 0.
  \label{eq:decomposition}
\end{equation}
\subsection{Mixed-Effects Model}
\label{sec:methods-mixed}

We complemented this decomposition with the mixed-effects model in equation~\eqref{eq:mixed}, fitted to the same 12,125 aggregate means using statsmodels (Seabold and Perktold). Here $\mu$ is the intercept, $\beta_f$ and $\tau_y$ are fixed field and year effects, $a_u$ is a random university intercept, and $b_{uf}$ is the random effect of a field within a university. We assumed independent, zero-mean normal random effects with variances $\sigma_a^2$ and $\sigma_b^2$, and normal residuals $\varepsilon_{ufy}$ with variance $\sigma^2$; $\mathcal N$ denotes the normal distribution. We estimated these variances by restricted maximum likelihood (REML) using the limited-memory Broyden--Fletcher--Goldfarb--Shanno optimiser (L-BFGS). Partial pooling draws less precisely estimated university profiles towards the shared pattern instead of treating every observed difference as equally reliable. We weighted aggregate means equally. For the descriptive comparison, we combined the spread of fitted fixed field effects with the interaction and residual variances, excluding university intercepts and year effects. This and the nested least-squares split estimate different quantities, so we compared their broad patterns rather than treating their percentages as interchangeable.
\begin{equation}
  \begin{aligned}
  \bar g_{ufy} &= \mu + \beta_f + \tau_y + a_u + b_{uf} + \varepsilon_{ufy},\\
  a_u &\sim \mathcal N(0,\sigma_a^2),\quad b_{uf} \sim \mathcal N(0,\sigma_b^2),\quad \varepsilon_{ufy} \sim \mathcal N(0,\sigma^2).
  \end{aligned}
  \label{eq:mixed}
\end{equation}
\subsection{Student Intake and Other Possible Explanations}
\label{sec:methods-confounders}

We examined whether observed intake and size differences could account for the estimated field effects. Within the public-university sample, we added standardised admission grades and the logarithm of graduate counts, expressing both relative to their university-year averages. Standardisation gives the admission coefficient a common unit of one standard deviation, while the logarithm limits the influence of very large fields. The final figure uses 8,043 records from 47 universities. We compared field effects before and after these controls on the same figure sample. We then added dropout in a separate sensitivity model: stricter grading could itself increase dropout, so adjustment might remove part of the association we wish to describe. We therefore avoided interpreting dropout as an established cause. In Portugal, we repeated the intake comparison with individual entry grades standardised within admission route and the background controls in our linear model, providing a check where intake is observed for individual students.
\subsection{Selection Simulation and the Normal Approximation}
\label{sec:methods-simulation}

We translated field grade differences into selection rates because a pooled ranking allocates opportunities to students rather than to field averages. Since SIIU publishes summaries, we approximated grades within each university-field-year by a normal distribution with mean $\bar g_{ufy}$ and standard deviation $\sigma_{ufy}=\mathrm{IQR}_{ufy}/1.349$, where the interquartile range (IQR) is the third quartile minus the first. The divisor is the normal distribution's IQR in standard-deviation units. In equation~\eqref{eq:selection}, $\Phi$ is the standard normal cumulative distribution function, $c_{uy}$ the common grade cut-off, $w_{ufy}$ the field's share of graduates among included fields, and $s_{ufy}$ its proportion above that cut-off. We solved for $c_{uy}$ so the weighted selection rate equalled 20\%. We required at least four reported fields per university-year, leaving 664 university-years from 79 universities; the simulation represents those available fields, not necessarily the entire institution.
\begin{equation}
  s_{ufy} = 1-\Phi\!\left(\frac{c_{uy}-\bar g_{ufy}}{\sigma_{ufy}}\right),\qquad
  \sum_{f} w_{ufy}\, s_{ufy} = 0.20 .
  \label{eq:selection}
\end{equation}
We checked the normal approximation in four ways. We compared each reported mean with its median in estimated standard-deviation units, calculated probability outside the valid 5--10 range, and compared observed maxima with the expected maximum and pointwise 95\% predictive interval for the reported number $n$ of graduates. Finally, we repeated selection with normal distributions truncated to $[5,10]$, using identical records, graduate weights and underlying location and scale. Truncation changes the resulting means and spreads; we did not refit them. These are consistency and sensitivity checks, not tests establishing individual normality: aggregate summaries conceal the full distribution, and the maximum calculation treats fitted parameters and graduate counts as known and assumes independent grades.
\subsection{Uncertainty}
\label{sec:methods-uncertainty}

We accounted for repeated observations when reporting uncertainty. For field-effect figures we clustered standard errors by university, allowing errors across fields and up to nine years within an institution to be correlated (Cameron and Miller). For prediction errors and selection gaps, we constructed 95\% intervals from 1,000 bootstrap samples of whole universities, retaining each sampled institution's observations together (Efron). We used student resampling for the Portuguese ranking comparison and Fisher's z transformation for correlations, with random seed 42 for resampling. Mixed-model profile intervals are conditional empirical-Bayes intervals: they hold fitted fixed effects and variance parameters constant and therefore omit uncertainty from estimating those parameters. They concern aggregate means, not prediction intervals for individual students.
